\section{Problemstellung und Zielsetzung}
\label{sec:problemstellung}

Die Verfahren des Standes der Technik gehen mit der Faserverstärkung unterschiedlich um. VDI 2736 Blatt 2 rechnet mit Nennspannung und Geometriefaktoren, kennt die Faserorientierung nicht und nennt Festigkeitswerte nur für wenige unverstärkte Werkstoffe (\secref{sec:vdi2736}). RIKOR und die FVA-Workbench erfassen die Lastverteilung unter Verformung, rechnen aber mit der Nenngeometrie und einem richtungsunabhängigen Werkstoff (\secref{sec:numerik_fussgeometrien}). Der vereinfachte Festigkeitsnachweis nach Stommel et al. schreibt für kurzfaserverstärkte Werkstoffe ein isotropes Modell vor und erfasst die Festigkeitsminderung quer zur Faser pauschal über die Werkstoffklasse (\secref{sec:stommel_stand}); seine Versagensdefinition ist der erste makroskopische Anriss, die der Laufversuche der Zahnbruch (\secref{sec:experiment_stand}). Die Fußrundungsvarianten mit elliptischer, Bézier- und bionischer Kontur sind überwiegend an Stahlrädern und unverstärkten Kunststoffen belegt. Für kurzfaserverstärkte Kunststoffzahnräder fehlen Aussagen zu Gültigkeit und Grenzen dieser Verfahren, und Laufversuche an der Prüfverzahnung kst-E liegen für PA46-GF30 nicht vor.

In der Bachelorarbeit des Verfassers wurde der Nachweis nach Stommel et al. rechnerisch auf Kunststoffzahnräder angewendet; die experimentelle Validierung für diese Anwendung blieb offen \cite[S.~24 und 71]{Kuh24}. Ziel dieser Arbeit ist deshalb, die Zahnfußtragfähigkeit kurzfaserverstärkter Kunststoffzahnräder am Beispiel der Paarung Stahl gegen PA46-GF30 rechnerisch und experimentell zu untersuchen und daraus das Potenzial der Faserverstärkung zu beurteilen. Die Zahnfußspannung wird nach VDI 2736 Blatt 2, mit der FVA-Workbench, mit RIKOR und mit einem eigenen FE-Modell berechnet, das die Faserorientierung der Spritzgießsimulation in vier Werkstoffstufen vom isotropen bis zum anisotropen Modell übernimmt. Der Nachweis nach Stommel et al. wird wie veröffentlicht und mit einer eigenen Erweiterung auf die anisotropen Stufen geführt, S-Life Plastics dient als Gegenprobe. Die drei Fußgeometrien werden in dieselbe Kette eingesetzt. Am FZG-Kleinverspannungsprüfstand mit dem Achsabstand $52$\,mm werden Laufversuche bei drei Drehmomenten gefahren, die Prüfräder vor und nach dem Versuch vermessen. Aus dem Vergleich von Berechnung und Versuch werden die Verfahren nach Genauigkeit und Anwendbarkeit bewertet. Den Aufbau der Arbeit nennt \secref{sec:aufbau}.

% Herkunft:
% Absatz 1 (Lücke): 10_report/aufgabenstellung.tex Absatz 1; die Einordnung der Verfahren verweist auf 2.5 und 2.6
%   (Entwürfe kap2_5, kap2_6), deren Quellen dort stehen: VDI 2736 Blatt 2 ohne Orientierung und ohne sigma_FlimN
%   für PA46 (VDI 2736 Blatt 2 S. 14), RIKOR/Workbench Nenngeometrie und richtungsunabhängiger Werkstoff (FVA 30 XI,
%   FVA 892 I), Stommel Isotropie-Vorschrift (Sto18 S. 343 bis 346) und erster makroskopischer Anriss (S. 204),
%   Fußrundungsvarianten (Frühe 2012, Roth/Opferkuch 2017, Patente, Kassem 2022/2023), keine kst-E-Laufversuche mit
%   PA46-GF30 (Arbeitsmappe der Laufversuche).
% Absatz 2 (Ziel): BA-0486.txt Druckseite 24 (2.4) und 71 (6: Forschungsbedarf experimentelle Validierung);
%   aufgabenstellung.tex Absatz 2 und Arbeitspakete; vier Werkstoffstufen und Erweiterung: kap3_6, kap3_7;
%   drei Drehmomente: Arbeitsmappe. Nutzer 2026-10-18: Problemstellung und Zielsetzung hier nach dem Stand der
%   Technik, Aufbau in Kapitel 1, kurz und bündig, Vorgeschichte aus der Bachelorarbeit nennen.
% Bibwunsch: Kuh24 (litbib). Nomenklatur: keine. Offen: keine.
